% Printed142; principalPDF149: blank verso.
\par\clearpage\null\thispagestyle{empty}\clearpage
% Printed143; principalPDF150.
\begin{center}
{\Large\textbf{NACHLASS.}}\par\medskip
{\large\textbf{[KONFORME ABBILDUNG DES SPHÄROIDS\\IN DER EBENE.]}}\par\smallskip
\rule{18mm}{.4pt}\par\medskip
[1.]\par\smallskip
[Berechnung der geografischen Breite und Länge aus den ebenen rechtwinkligen Koordinaten.]
\end{center}
[Bei der konformen Übertragung der Sphäroidfläche auf die Ebene wird ein Meridian, der Hauptmeridian, durch eine Gerade, die \(x\)-Achse, dargestellt. Jeder Abschnitt auf der Abszissenachse ist dem Teil des Hauptmeridians gleich, dessen Bild er ist.

Es sei \(a\) die große Halbachse und \(e\) die Exzentrizität der Meridianellipse. Einem Punkt auf der Ellipsoidfläche, dessen Breite \(\Phi\) und dessen Länge \(\lambda\) ist, sollen in der Ebene die rechtwinkligen Koordinaten \(x,y\) entsprechen. Zum Schnittpunkt des Parallelkreises von der Breite \(\Phi\) mit dem Hauptmeridian gehöre die Abszisse \(\xi\), und zum Endpunkt der Abszisse \(x\) die Breite \(\varphi\). \(\xi\) und \(x\) sind also zugleich die entsprechenden Meridianbögen, vom Äquator an gemessen.]

Es sei ein Linienelement
\[=\sqrt{\{\mathrm d\Phi^2+(\Theta(\Phi))^2\mathrm d\lambda^2\}};\]
\noindent man setze
\[\int\frac{\mathrm d\Phi}{\Theta(\Phi)}=\mathrm f(\xi),\]

% Printed144; principalPDF151.
\par\clearpage
\noindent144\hfill{\small\textbf{NACHLASS.}}\hfill\null\par\medskip
\noindent so wird
\[\mathrm f(x+iy)=\mathrm f(\xi)+i\lambda=\mathrm F(\Phi)+i\lambda.\]
[Für das Sphäroid ist
\[
\begin{aligned}
\mathrm ds^2&=\frac{aa(1-ee)^2}{(1-ee\sin\Phi^2)^3}\mathrm d\Phi^2+\frac{aa\cos\Phi^2}{1-ee\sin\Phi^2}\mathrm d\lambda^2\\
&=\frac{aa(1-ee)^2}{(1-ee\sin\Phi^2)^3}\GaussOldenburgBracePair{\mathrm d\Phi^2+\left(\frac{(1-ee\sin\Phi^2)\cos\Phi}{1-ee}\right)^2\mathrm d\lambda^2}.
\end{aligned}
\]
Entwickelt man \(\mathrm f(x+iy)\) nach Potenzen von \(iy\), so ergibt sich:
\[
\begin{aligned}
\mathrm f(\xi)=\mathrm F(\Phi)&=\mathrm f(x)-\tfrac12yy\mathrm f''(x)+\tfrac1{24}y^4\mathrm f^{\mathrm{IV}}(x)-\ldots\\
\lambda&=y\mathrm f'(x)-\tfrac16y^3\mathrm f'''(x)+\tfrac1{120}y^5\mathrm f^{\mathrm V}(x)-\ldots.
\end{aligned}
\]
Für \(y=0\) ist \(\Phi=\varphi\), also]
\[\mathrm F(\varphi)=\mathrm f(x).\]
[Da \(\dfrac1{\Theta(\varphi)}=\mathrm F'(\varphi)\), so ist]
\[\mathrm F'(\varphi)=\frac{1-ee}{(1-ee\sin\varphi^2)\cos\varphi}.\]
[\(x\) wächst mit \(\varphi\), also]
\[\mathrm dx=\frac{a(1-ee)\mathrm d\varphi}{(1-ee\sin\varphi^2)^{\frac32}},\]
\noindent[mithin wird]
\[\mathrm f'(x)=\frac{(1-ee\sin\varphi^2)^{\frac12}}{a\cos\varphi},\]
\noindent[und weiter]
\[
\begin{aligned}
\mathrm f''(x)&=\frac{(1-ee\sin\varphi^2)\sin\varphi}{aa\cos\varphi^2}\\
\mathrm f'''(x)&=\frac{(1-ee\sin\varphi^2)^{\frac32}}{a^3(1-ee)\cos\varphi^3}\{1+(1-3ee)\sin\varphi^2+ee\sin\varphi^4\}\\
&\left[=\frac{(1-ee\sin\varphi^2)^{\frac32}}{a^3(1-ee)\cos\varphi^3}\{(1-ee)(2-\cos\varphi^2)+ee\cos\varphi^4\}\right]\\
\mathrm f^{\mathrm{IV}}(x)&=\frac{(1-ee\sin\varphi^2)^2\sin\varphi}{a^4(1-ee)^2\cos\varphi^4}\{5-9ee+(1-4ee+15e^4)\sin\varphi^2\\
&\hspace{32mm}+(ee-13e^4)\sin\varphi^4+4e^4\sin\varphi^6\}\\
&\left[=\frac{(1-ee\sin\varphi^2)^2\sin\varphi}{a^4(1-ee)^2\cos\varphi^4}\{(1-ee)^2(6-\cos\varphi^2)+ee(1-ee)\cos\varphi^4-4e^4\cos\varphi^6\}\right]
\end{aligned}
\]
{\centering usw.\par}
[Das Gesetz der Differentiation ist das folgende.]

% Printed145; principalPDF152.
\par\clearpage
\noindent\null\hfill{\small\textbf{KONFORME ABBILDUNG DES SPHÄROIDS IN DER EBENE.}}\hfill145\par\medskip
Man setze
\[\frac{\sqrt{(1-ee\sin\varphi^2)}}{a\cos\varphi}=t,\qquad[\sin\varphi=s]\]
\noindent und
\[\mathrm f^n(x)=t^nu,\]
\noindent so ist
\[\mathrm f^{n+1}(x)=t^{n+1}\left\{nus+\frac{\partial u}{\partial s}.\frac{(1-ss)(1-eess)}{1-ee}\right\}\]
\noindent oder, wenn man die Potenzen von \(ee\) vernachlässigt,
\[\mathrm f^{n+1}(x)=t^{n+1}\left\{nus+\frac{\partial u}{\partial s}(1-ss)(1+ee(1-ss))\right\}.\]
So findet man für
\[
\begin{array}{c|c}
n&u\\\hline
1&1\\2&s\\3&1+ss+ee(1-ss)^2\\4&5s+s^3+ee(1-ss)^2s
\end{array}
\]
[Die Substitution der Werte von \(\mathrm f'(x)\) und \(\mathrm f'''(x)\) in der Gleichung für \(\lambda\) gibt:]
\[
\lambda=\frac{(1-ee\sin\varphi^2)^{\frac12}}{a\cos\varphi}y-\frac{(1-ee\sin\varphi^2)^{\frac32}}{6a^3(1-ee)\cos\varphi^3}(2(1-ee)-(1-ee)\cos\varphi^2+ee\cos\varphi^4)y^3\ldots
\]
[Setzt man]
\[\mathrm F(\Phi)=\mathrm F(\varphi)+\omega,\]
\noindent[wo also
\[\omega=-\tfrac12yy\mathrm f''(x)+\tfrac1{24}y^4\mathrm f^{\mathrm{IV}}(x)\ldots\]
\noindent ist, so wird]
\[\begin{aligned}\Phi&=\mathrm G\{\mathrm F(\varphi)+\omega\}\\&=\varphi+A\omega+B\omega\omega+C\omega^3+\ldots,\end{aligned}\]
\noindent[wobei]
\[A=\left[\frac{\mathrm d\varphi}{\mathrm dF}=\right]\frac{(1-ee\sin\varphi^2)\cos\varphi}{1-ee}.\]
[Wenn
\[-\mathrm F'(\varphi)=\frac{\partial\omega}{\partial\varphi}\]
\noindent gesetzt wird, so ist aber]
\[
\begin{aligned}
&1+\frac{\partial A}{\partial\varphi}.\omega+\frac{\partial B}{\partial\varphi}.\omega\omega+\frac{\partial C}{\partial\varphi}.\omega^3\ldots\\
&=A\mathrm F'+2B\mathrm F'.\omega+[3C\mathrm F'.\omega\omega+4D\mathrm F'.\omega^3+]\ldots;
\end{aligned}
\]
\noindent{\small IX.\hfill19}\par

% Printed146; principalPDF153.
\par\clearpage
\noindent146\hfill{\small\textbf{NACHLASS.}}\hfill\null\par\medskip
\noindent[daraus folgt:]
\[
\begin{array}{cc}
A\mathrm F'=1&2B=AA'\\2B\mathrm F'=A'&3C=AB'\\3C\mathrm F'=B'&4D=AC'\\\text{etc.}&\text{etc.}
\end{array}
\]
\[
\begin{aligned}
A'&=-\frac{\sin\varphi}{1-ee}(1+2ee-3ee\sin\varphi^2)\\
B&=-\frac{1-ee\sin\varphi^2}{2(1-ee)^2}\sin\varphi\cos\varphi(1+2ee-3ee\sin\varphi^2).
\end{aligned}
\]
[Wird in der Gleichung \(\Phi=\varphi+A\omega+B\omega\omega+\ldots\) für \(\omega\) der Wert eingesetzt, so ist zunächst
\[\Phi=\varphi-\tfrac12A\mathrm f''(x).yy+\left(\tfrac1{24}A\mathrm f^{\mathrm{IV}}(x)+\tfrac14B(\mathrm f''(x))^2\right)y^4\ldots;\]
\noindent mit den obigen Werten für \(A\) und \(B\) und den Werten von \(\mathrm f''(x)\) und \(\mathrm f^{\mathrm{IV}}(x)\), S.~144, erhält man daher:]
\[
\begin{aligned}
\Phi={}&\varphi-\frac{(1-ee\sin\varphi^2)^2\operatorname{tang}\varphi}{2aa(1-ee)}yy\\
&+\frac{(1-ee\sin\varphi^2)^3\sin\varphi}{24a^4(1-ee)^3\cos\varphi^3}\{5-9ee-(2+7ee-21e^4)\sin\varphi^2\\
&\hspace{45mm}+(10ee-22e^4)\sin\varphi^4+4e^4\sin\varphi^6\}y^4\ldots
\end{aligned}
\]
\noindent[oder]
\[
\begin{aligned}
\Phi={}&\varphi-\frac{(1-ee\sin\varphi^2)^2\operatorname{tang}\varphi}{2aa(1-ee)}yy\\
&+\frac{(1-ee\sin\varphi^2)^3\sin\varphi}{24a^4(1-ee)^3\cos\varphi^3}\{3(1-ee)^2+(2-11ee)(1-ee)\cos\varphi^2\\
&\hspace{45mm}+10ee(1-ee)\cos\varphi^4-4e^4\cos\varphi^6\}y^4\ldots
\end{aligned}
\]
[Wenn \(\Phi=\varphi-z\) gesetzt wird, so ist
\[\log\operatorname{hyp}\sin\Phi=\log\operatorname{hyp}\sin\varphi-\frac1{\operatorname{tang}\varphi}z-\tfrac12\frac1{\sin\varphi^2}zz\ldots;\]
\noindent mithin:]
\[
\begin{aligned}
\log\operatorname{hyp}\sin\Phi={}&\log\operatorname{hyp}\sin\varphi-\frac{(1-ee\sin\varphi^2)^2}{2aa(1-ee)}yy\\
&+\frac{(1-ee\sin\varphi^2)^3}{12a^4(1-ee)^2}\left\{1-7ee+5ee\cos\varphi^2-\frac{2e^4}{1-ee}\cos\varphi^4\right\}y^4\ldots.
\end{aligned}
\]
{\centering [2.]\par
[Berechnung der Meridiankonvergenz aus den ebenen rechtwinkligen Koordinaten.]\par}
[Die Meridiankonvergenz \(c\) ist der Winkel, den das Bild des Meridians mit der Parallelen zur \(x\)-Achse bildet. \(x\) sei nach Norden positiv, dann ist:
\[\operatorname{tang}c=-\frac{\mathrm dy}{\mathrm dx}.\]

% Printed147; principalPDF154.
\par\clearpage
\noindent\null\hfill{\small\textbf{KONFORME ABBILDUNG DES SPHÄROIDS IN DER EBENE.}}\hfill147\par\medskip
Aus der Gleichung, die in der Ebene den Meridian von der Länge \(\lambda\) darstellt:
\[\lambda=\operatorname{const.}=y\mathrm f'(x)-\tfrac16y^3\mathrm f'''(x)\ldots,\]
\noindent folgt aber
\[\frac{\mathrm dy}{\mathrm dx}=-\frac{y\mathrm f''(x)-\tfrac16y^3\mathrm f^{\mathrm{IV}}(x)\ldots}{\mathrm f'(x)-\tfrac12yy\mathrm f'''(x)\ldots},\]
\noindent also ist:]
\[
\begin{aligned}
\operatorname{tang}c&=\frac{y\mathrm f''(x)-\tfrac16y^3\mathrm f^{\mathrm{IV}}(x)\ldots}{\mathrm f'(x)-\tfrac12yy\mathrm f'''(x)\ldots}\\
&\left[=y\frac{\mathrm f''(x)}{\mathrm f'(x)}\left\{1-yy\left(\tfrac16\frac{\mathrm f^{\mathrm{IV}}(x)}{\mathrm f''(x)}-\tfrac12\frac{\mathrm f'''(x)}{\mathrm f'(x)}\right)\ldots\right\}. \right.
\end{aligned}
\]
Daher wird nach S.~144]
\[\operatorname{tang}c=D.\frac{\sqrt{(1-ee\sin\varphi^2)}}a y\operatorname{tang}\varphi,\]
\noindent wo
\[
\begin{aligned}
\log\operatorname{hyp}D&=\left(-\tfrac16\frac{\mathrm f^{\mathrm{IV}}(x)}{\mathrm f''(x)}+\tfrac12\frac{\mathrm f'''(x)}{\mathrm f'(x)}\right)yy\ldots\\
&=-\tfrac13yy\frac{(1-ee\sin\varphi^2)^2}{aa(1-ee)^2}(1-3ee+2ee\sin\varphi^2)\ldots
\end{aligned}
\]
\noindent[oder näherungsweise]
\[=-\tfrac13\frac{yy}{aa}(1-ee).\]
[Ferner hat man:]
\[c=y\frac{\mathrm f''(x)}{\mathrm f'(x)}-\tfrac16y^3\left\{\frac{\mathrm f^{\mathrm{IV}}(x)}{\mathrm f'(x)}-\frac{3\mathrm f''(x)\mathrm f'''(x)}{(\mathrm f'(x))^2}+2\left(\frac{\mathrm f''(x)}{\mathrm f'(x)}\right)^3\right\}\ldots\]
\noindent[oder, wenn
\[\frac{\mathrm f''(x)}{\mathrm f'(x)}=\mathrm h(x)\]
\noindent gesetzt wird,
\[c=y\mathrm h(x)-\tfrac16y^3\mathrm h''(x)\ldots].\]
\[
\begin{aligned}
\mathrm h(x)&=\frac{(1-ee\sin\varphi^2)^{\frac12}\sin\varphi}{a\cos\varphi}\\
\mathrm h'(x)&=\frac{1-ee\sin\varphi^2}{aa(1-ee)\cos\varphi^2}\{1-2ee\sin\varphi^2+ee\sin\varphi^4\}\\
{[}\mathrm h''(x)&=\frac{2(1-ee\sin\varphi^2)^{\frac32}\sin\varphi}{a^3(1-ee)^2\cos\varphi^3}\{1-3ee+(2ee+4e^4)\sin\varphi^2\\
&\hspace{37mm}-(ee+5e^4)\sin\varphi^4+2e^4\sin\varphi^6\}\\
&=\frac{2(1-ee\sin\varphi^2)^{\frac32}\sin\varphi}{a^3(1-ee)^2\cos\varphi^3}\{(1-ee)^2-ee(1-ee)\cos\varphi^4-2e^4\cos\varphi^6\}].
\end{aligned}
\]
\noindent\hfill{\small19*}\par

% Printed148; principalPDF155.
\par\clearpage
\noindent148\hfill{\small\textbf{NACHLASS.}}\hfill\null\par\medskip
Also ist:
\[
\begin{aligned}
c={}&\frac{\sqrt{(1-ee\sin\varphi^2)}}a\operatorname{tang}\varphi.y\\
&-\tfrac13\frac{(1-ee\sin\varphi^2)^{\frac32}\sin\varphi}{a^3(1-ee)^2\cos\varphi^3}\{1-3ee+(2ee+4e^4)\sin\varphi^2\\
&\hspace{44mm}-(ee+5e^4)\sin\varphi^4+2e^4\sin\varphi^6\}y^3\ldots
\end{aligned}
\]
\noindent[oder]
\[
\begin{aligned}
c={}&\frac{\sqrt{(1-ee\sin\varphi^2)}}a\operatorname{tang}\varphi.y\\
&-\tfrac13\frac{(1-ee\sin\varphi^2)^{\frac32}\sin\varphi}{a^3(1-ee)^2\cos\varphi^3}\{(1-ee)^2\quad *\quad-ee(1-ee)\cos\varphi^4-2e^4\cos\varphi^6\}y^3\ldots.
\end{aligned}
\]
[Man hat auch, wenn man \(\lambda\) einführt:]
\[c=\lambda\sin\varphi-\frac{(1-ee\sin\varphi^2)^{\frac32}\operatorname{tang}\varphi}{6a^3(1-ee)^2}\{(1-ee)^2-3ee(1-ee)\cos\varphi^2-4e^4\cos\varphi^4\}y^3\ldots.\]
\medskip
{\centering [3.]\par}
Zur numerischen Berechnung sind folgende Formeln am bequemsten. Es sei
\[
\begin{aligned}
L&=\frac{(1-ee\sin\varphi^2)^{\frac12}y}{a\rho\cos\varphi}\\
M&=\frac{(1-ee\sin\varphi^2)^2\operatorname{tang}\varphi.yy}{2aa(1-ee)\rho}\\
N&=\frac{(1-ee\sin\varphi^2)^{\frac12}}{a\rho}\operatorname{tang}\varphi.y.
\end{aligned}
\]
Dann ist
\[
\begin{aligned}
\log\lambda&=\log L-A.LL\\
\log(\varphi-\Phi)&=\log M-B.LL\\
\log c&=\log N-C.LL,
\end{aligned}
\]
\noindent wo die Logarithmen von \(A,B,C\) am bequemsten aus einer besonderen Tafel mit dem Argument \(\varphi\) genommen werden.

Zur Berechnung dieser Tafel dienen die Formeln:
\[
\begin{aligned}
A&=H\left\{0{,}75-\tfrac14\cos2\varphi+\frac{ee}{2(1-ee)}\cos\varphi^4\right\}\\
B&=H\left\{0{,}75+\frac{2-11ee}{4(1-ee)}\cos\varphi^2+\frac{5ee}{2(1-ee)}\cos\varphi^4-\frac{e^4}{(1-ee)^2}\cos\varphi^6\right\}\\
C&=H\left\{1-\frac{ee}{1-ee}\cos\varphi^4-\frac{2e^4}{(1-ee)^2}\cos\varphi^6\right\},
\end{aligned}
\]

% Printed149; principalPDF156.
\par\clearpage
\noindent\null\hfill{\small\textbf{KONFORME ABBILDUNG DES SPHÄROIDS IN DER EBENE.}}\hfill149\par\medskip
\noindent wo
\[H=\tfrac13k\rho\rho.10^7,\qquad\log H=5{,}531\,8128-10\]
[\(k=\) Modul der Briggs'schen Logarithmen und \(\rho=\dfrac1{206264{,}8\ldots}\)

Die Korrekturen \(A.LL,B.LL,C.LL\) werden in Einheiten der 7.~Dezimalstelle erhalten. \(\lambda,\varphi-\Phi\) und \(c\) erhält man in Bogensekunden.

Der nachfolgenden Tabelle für die Logarithmen von \(A,B,C\) liegt der Abplattungswert \(\frac1{302{,}68}\) zugrunde].
\begin{center}
\begingroup\scriptsize\setlength{\tabcolsep}{2pt}\renewcommand{\arraystretch}{1.15}
\begin{tabular}[t]{|r|r|r|r|}\hline
\(\varphi\)&\(\log A\)&\(\log B\)&\(\log C\)\\\hline\hline
&\(-10\)&\(-10\)&\(-10\)\\
\(50^\circ\;0'\)&\(5{,}431\,6218\)&\(5{,}510\,4149\)&\(5{,}531\,3175\)\\
\(3\)&\(8555\)&\(2206\)&\(95\)\\
\(6\)&\(5{,}432\,0889\)&\(0263\)&\(5{,}531\,3216\)\\
\(9\)&\(3222\)&\(5{,}509\,8319\)&\(36\)\\
\(12\)&\(5552\)&\(6376\)&\(57\)\\
\(15\)&\(7881\)&\(4432\)&\(77\)\\
\hline
\(18\)&\(5{,}433\,0208\)&\(2488\)&\(97\)\\
\(21\)&\(2532\)&\(0544\)&\(5{,}531\,3318\)\\
\(24\)&\(4854\)&\(5{,}508\,8600\)&\(38\)\\
\(27\)&\(7174\)&\(6656\)&\(59\)\\
\(30\)&\(9494\)&\(4712\)&\(79\)\\
\hline
\(33\)&\(5{,}434\,1810\)&\(2768\)&\(99\)\\
\(36\)&\(4124\)&\(0823\)&\(5{,}531\,3419\)\\
\(39\)&\(6438\)&\(5{,}507\,8878\)&\(39\)\\
\(42\)&\(8748\)&\(6934\)&\(59\)\\
\(45\)&\(5{,}435\,1057\)&\(4989\)&\(79\)\\
\hline
\(48\)&\(3364\)&\(3044\)&\(99\)\\
\(51\)&\(5668\)&\(1100\)&\(5{,}531\,3519\)\\
\(54\)&\(7971\)&\(5{,}506\,9155\)&\(38\)\\
\(57\)&\(5{,}436\,0271\)&\(7210\)&\(58\)\\
\(51\quad0\)&\(2570\)&\(5265\)&\(78\)\\
\hline
\(3\)&\(4867\)&\(3320\)&\(98\)\\
\(6\)&\(7161\)&\(1375\)&\(5{,}531\,3617\)\\
\(9\)&\(9454\)&\(5{,}505\,9430\)&\(37\)\\
\(12\)&\(5{,}437\,1745\)&\(7485\)&\(56\)\\
\(15\)&\(4033\)&\(5540\)&\(76\)\\
\hline
\(18\)&\(6319\)&\(3595\)&\(95\)\\
\(21\)&\(8604\)&\(1650\)&\(5{,}531\,3714\)\\
\(24\)&\(5{,}438\,0886\)&\(5{,}504\,9704\)&\(34\)\\
\(27\)&\(3166\)&\(7759\)&\(53\)\\
\(30\)&\(5444\)&\(5814\)&\(72\)\\
\hline
\(33\)&\(7720\)&\(3869\)&\(91\)\\
\(36\)&\(9994\)&\(1924\)&\(5{,}531\,3810\)\\
\(39\)&\(5{,}439\,2266\)&\(5{,}503\,9979\)&\(30\)\\
\(42\)&\(4535\)&\(8034\)&\(49\)\\
\(45\)&\(6803\)&\(6089\)&\(68\)\\
\hline\end{tabular}
\hfill
\begin{tabular}[t]{|r|r|r|r|}\hline
\(\varphi\)&\(\log A\)&\(\log B\)&\(\log C\)\\\hline\hline
&\(-10\)&\(-10\)&\(-10\)\\
\(51^\circ45'\)&\(5{,}439\,6803\)&\(5{,}503\,6089\)&\(5{,}531\,3868\)\\
\(48\)&\(9068\)&\(4144\)&\(87\)\\
\(51\)&\(5{,}440\,1332\)&\(2199\)&\(5{,}531\,3906\)\\
\(54\)&\(3593\)&\(0254\)&\(24\)\\
\(57\)&\(5852\)&\(5{,}502\,8310\)&\(43\)\\
\(52\quad0\)&\(8109\)&\(6365\)&\(62\)\\
\hline
\(3\)&\(5{,}441\,0364\)&\(4420\)&\(81\)\\
\(6\)&\(2616\)&\(2476\)&\(99\)\\
\(9\)&\(4867\)&\(0531\)&\(5{,}531\,4018\)\\
\(12\)&\(7115\)&\(5{,}501\,8587\)&\(36\)\\
\(15\)&\(9361\)&\(6643\)&\(55\)\\
\hline
\(18\)&\(5{,}442\,1605\)&\(4699\)&\(73\)\\
\(21\)&\(3847\)&\(2755\)&\(91\)\\
\(24\)&\(6087\)&\(0812\)&\(5{,}531\,4110\)\\
\(27\)&\(8325\)&\(5{,}500\,8868\)&\(28\)\\
\(30\)&\(5{,}443\,0560\)&\(6925\)&\(46\)\\
\hline
\(33\)&\(2793\)&\(4982\)&\(64\)\\
\(36\)&\(5024\)&\(3039\)&\(82\)\\
\(39\)&\(7254\)&\(1096\)&\(5{,}531\,4200\)\\
\(42\)&\(9480\)&\(5{,}499\,9153\)&\(18\)\\
\(45\)&\(5{,}444\,1705\)&\(7211\)&\(36\)\\
\hline
\(48\)&\(3927\)&\(5269\)&\(54\)\\
\(51\)&\(6148\)&\(3327\)&\(72\)\\
\(54\)&\(8366\)&\(1384\)&\(89\)\\
\(57\)&\(5{,}445\,0581\)&\(5{,}498\,9443\)&\(5{,}531\,4307\)\\
\(53\quad0\)&\(2795\)&\(7501\)&\(25\)\\
\hline
\(3\)&\(5006\)&\(5560\)&\(43\)\\
\(6\)&\(7215\)&\(3618\)&\(60\)\\
\(9\)&\(9422\)&\(1677\)&\(78\)\\
\(12\)&\(5{,}446\,1627\)&\(5{,}497\,9737\)&\(95\)\\
\(15\)&\(3829\)&\(7796\)&\(5{,}531\,4413\)\\
\hline
\(18\)&\(6029\)&\(5856\)&\(30\)\\
\(21\)&\(8227\)&\(3916\)&\(48\)\\
\(24\)&\(5{,}447\,0423\)&\(1976\)&\(65\)\\
\(27\)&\(2617\)&\(0036\)&\(83\)\\
\(30\)&\(4808\)&\(5{,}496\,8097\)&\(5{,}531\,4500\)\\
\hline\end{tabular}
\hfill
\begin{tabular}[t]{|r|r|r|r|}\hline
\(\varphi\)&\(\log A\)&\(\log B\)&\(\log C\)\\\hline\hline
&\(-10\)&\(-10\)&\(-10\)\\
\(53^\circ30'\)&\(5{,}447\,4808\)&\(5{,}496\,8097\)&\(5{,}531\,4500\)\\
\(33\)&\(6997\)&\(6158\)&\(17\)\\
\(36\)&\(9184\)&\(4219\)&\(34\)\\
\(39\)&\(5{,}448\,1369\)&\(2281\)&\(51\)\\
\(42\)&\(3551\)&\(0343\)&\(68\)\\
\(45\)&\(5731\)&\(5{,}495\,8405\)&\(85\)\\
\hline
\(48\)&\(7909\)&\(6468\)&\(5{,}531\,4602\)\\
\(51\)&\(5{,}449\,0084\)&\(4530\)&\(19\)\\
\(54\)&\(2258\)&\(2594\)&\(35\)\\
\(57\)&\(4428\)&\(0657\)&\(52\)\\
\(54\quad0\)&\(6597\)&\(5{,}494\,8721\)&\(5{,}531\,4669\)\\
\hline
\(3\)&\(8763\)&\(6785\)&\(86\)\\
\(6\)&\(5{,}450\,0927\)&\(4850\)&\(5{,}531\,4702\)\\
\(9\)&\(3089\)&\(2915\)&\(19\)\\
\(12\)&\(5248\)&\(0980\)&\(35\)\\
\(15\)&\(7405\)&\(5{,}493\,9046\)&\(52\)\\
\hline
\(18\)&\(9560\)&\(7112\)&\(68\)\\
\(21\)&\(5{,}451\,1712\)&\(5179\)&\(84\)\\
\(24\)&\(3863\)&\(3245\)&\(5{,}531\,4801\)\\
\(27\)&\(6010\)&\(1313\)&\(17\)\\
\(30\)&\(8156\)&\(5{,}492\,9380\)&\(33\)\\
\hline
\(33\)&\(5{,}452\,0299\)&\(7448\)&\(49\)\\
\(36\)&\(2440\)&\(5516\)&\(65\)\\
\(39\)&\(4579\)&\(3585\)&\(81\)\\
\(42\)&\(6715\)&\(1654\)&\(97\)\\
\(45\)&\(8849\)&\(5{,}491\,9724\)&\(5{,}531\,4913\)\\
\hline
\(48\)&\(5{,}453\,0981\)&\(7794\)&\(29\)\\
\(51\)&\(3110\)&\(5865\)&\(45\)\\
\(54\)&\(5237\)&\(3936\)&\(60\)\\
\(57\)&\(7362\)&\(2007\)&\(76\)\\
\(55\quad0\)&\(9484\)&\(0079\)&\(92\)\\
\hline\end{tabular}
\endgroup
\end{center}

% Printed150; principalPDF157.
\par\clearpage
\noindent150\hfill{\small\textbf{NACHLASS.}}\hfill\null\par\medskip
{\centering [4.]\par}
[Es ist auch]
\[\lambda=A'\frac{\sqrt{(1-ee\sin\varphi^2)}}a.\frac y{\cos\varphi},\]
\noindent wo
\[\log\operatorname{hyp}A'=-\tfrac13cc-\tfrac16\frac{(1-ee\sin\varphi^2)^2}{aa(1-ee)}yy\ldots\]
\noindent oder
\[=-\tfrac13\lambda\lambda+\frac{1-ee}{6aa}\left(1-\frac{e^4}{(1-ee)^2}\cos\varphi^4\right)yy\ldots\]
\noindent[und näherungsweise]
\[=-\tfrac13\lambda\lambda+\frac{1-ee}{6aa}yy.\]
[Ferner ist]
\[\Phi=\varphi-B'\frac{(1-ee\sin\varphi^2)^2}{2aa(1-ee)}\operatorname{tang}\varphi.yy,\]
\noindent wo [näherungsweise]
\[
\begin{aligned}
\log\operatorname{hyp}B'&=-\tfrac14\lambda\lambda-\tfrac16\frac{yy}{aa}\\
&=-\tfrac5{12}\lambda\lambda+\tfrac16cc\\
&=-\tfrac14cc-\tfrac5{12}\frac{yy}{aa}.
\end{aligned}
\]
Auch kann man setzen
\[c=C'\frac{\sqrt{(1-ee\sin\varphi^2)}}a\operatorname{tang}\varphi.y,\]
\noindent wo
\[
\log\operatorname{hyp}C'=-\tfrac13cc-\tfrac13\left\{\frac{1-ee\sin\varphi^2}{aa}-\frac{ee(1-ee\sin\varphi^2)\cos\varphi^2(1-ee+2ee\cos\varphi^2)}{aa(1-ee)^2}\right\}yy\ldots
\]
\noindent[oder näherungsweise
\[=-\tfrac13cc-\tfrac13\frac{1-ee}{aa}yy].\]

% Printed151; principalPDF158.
\par\clearpage
\noindent\null\hfill{\small\textbf{KONFORME ABBILDUNG DES SPHÄROIDS IN DER EBENE.}}\hfill151\par\medskip
{\centering [5.]\par}
Wenn die Lage der Punkte durch Koordinaten \(x,y\) dargestellt wird, wobei die erste Koordinate im Meridian gleichförmig wächst, findet man die Breite \(\Phi\) des Orts, die Konvergenz seines Meridians \(c\) und seine Länge \(\lambda\) auf folgende Art.

Es sei \(\varphi\) die zur Abszisse \(x\) im Meridian gehörige Breite,
\[
\begin{aligned}
q&=\operatorname{Kompl.}\log\sqrt{(1-ee\sin\varphi^2)}\\
E&=\frac{1-ee}{6aa}.k.10^7\\
H&=\frac{k.10^7}{3.206265^2}.
\end{aligned}
\]
[\(k\) Modul der Briggs'schen Logarithmen
\[
\begin{aligned}
\log B&=\log\frac{206264{,}8\ldotp\ldotp}a-q,\qquad\log A=\log\frac{206264{,}8\ldotp\ldotp}{a(1-ee)}-3q\\
\log C&=\log\frac{206264{,}8\ldotp\ldotp}{2aa(1-ee)}-4q.
\end{aligned}
\]
Für den Abplattungswert \(\frac1{302{,}68}\) ist:]
\[
\begin{aligned}
\log E&=2{,}24756\,29768-10,\qquad\log H=5{,}53181\,27903-10\\
\log B&=8{,}50982\,72984-10-q\\
\log C&=\log\frac{AB}{2.206264{,}8\ldotp\ldotp}=1{,}40707\,38825-10-4q.
\end{aligned}
\]
[Zur Berechnung von \(A\) und \(B\) lässt sich direkt die Tabelle von S.~84 benutzen, oder man kann \(q\) aus der Tabelle von S.~77 entnehmen.]
\[
\begin{aligned}
c&=(1)By\operatorname{tang}\varphi&\log(1)&=-2Eyy-Hcc\\
\lambda&=(2)B\frac y{\cos\varphi}&\log(2)&=+Eyy-H\lambda\lambda\\
\Phi&=\varphi-(3)Cyy\operatorname{tang}\varphi&\log(3)&=-\tfrac34H\lambda\lambda-Eyy\\
&&&=-\tfrac54H\lambda\lambda+\tfrac12Hcc\\
&&&=-\tfrac34Hcc-\tfrac52Eyy.
\end{aligned}
\]
[\(\log(1),\log(2),\log(3)\) werden in Einheiten der 7.~Dezimalstelle erhalten. \(c,\lambda,\Phi-\varphi\) ergeben sich in Bogensekunden.]
